% !PRINTED 470
\weberpage{470}

%<<CONT
numerical factors, as follows easily from (8), are in order
\[
h,\ g,\ h^2,\ hg,\ (g^2,h^3),\ (g^2,h^2g),
\]
in which $(g^2,h^3)$, $(g^2,h^2g)$ denote linear homogeneous expressions
with numerical coefficients.

The calculation can be carried out just as in the special case
above. For simplification one may set $x_3=0$ and still obtains
equations enough for the determination of all the coefficients.
One finds in this way the coefficients in order:
\[
7\,\frac{1-\sqrt{-7}}{2}\,h,
\]
\[
-7\,\frac{1-\sqrt{-7}}{2}\,g,
\]
\[
-7\,(4+\sqrt{-7})\,h^2,
\]
\[
14\,(2+\sqrt{-7})\,hg,
\]
\[
-7\,\frac{5+\sqrt{-7}}{2}\,g^2 - 7\cdot\frac{7+3\sqrt{-7}}{2}\,h^3,
\]
\[
-g^2 + \frac{167-7\sqrt{-7}}{2}\,gh^2.
\]
The functions $g, h$ are fractional invariants, which depend only
on the ratios $x_1:x_2:x_3$, and we can express every other invariant
of the same property \gesperrt{rationally} in terms of $g, h$. For if
we represent such an invariant as a quotient of two forms of equal
degree without common divisor, then numerator and denominator
must be integral invariants of equal degree, because two fractional
functions of the variables $x$ represented in simplest form can be
equal to one another only when numerator and denominator are
separately equal to one another up to constant factors. Now here
numerator and denominator cannot be of odd degree, because
otherwise they would have the common factor $K$ (\S.\,126). Hence
numerator and denominator are representable rationally in terms of $f, C, \Delta$;
and if a term occurring in the numerator or in the denominator is
\weq{13}{f^a\,\Delta^b\,C^c,}
and $m$ is the degree of numerator and denominator, then
\weq{14}{4a + 6b + 14c = m.}
